% !TeX program = xelatex
\RequirePackage{fix-cm}
\documentclass[10pt,letterpaper,nopreprintline]{elsarticle}
\usepackage[letterpaper,left=17mm,right=17mm,top=28mm,bottom=28mm]{geometry}
\usepackage[OT1]{fontenc}
\renewcommand{\rmdefault}{cmr}
\renewcommand{\sfdefault}{cmss}
\renewcommand{\ttdefault}{cmtt}
\usepackage{setspace}
\setstretch{1.15}
\usepackage{amsmath,amssymb}
\usepackage{graphicx}
\usepackage{booktabs,array,longtable}
\usepackage{enumitem}
\usepackage{placeins}
\usepackage{needspace}
\usepackage[font=small,labelfont=bf,labelsep=period]{caption}
\usepackage[ruled,linesnumbered,longend]{algorithm2e}
\usepackage{xurl}
\usepackage{xcolor}
\definecolor{linkblue}{HTML}{1597D0}
\usepackage[colorlinks=true,linkcolor=linkblue,citecolor=linkblue,
    urlcolor=linkblue,filecolor=linkblue,anchorcolor=linkblue,
    breaklinks=true,bookmarksnumbered=true,unicode]{hyperref}
\AtBeginDocument{\hypersetup{linkcolor=linkblue,
    citecolor=linkblue,urlcolor=linkblue,filecolor=linkblue,
    anchorcolor=linkblue,menucolor=linkblue}}
\biboptions{numbers,sort&compress}
\setcounter{secnumdepth}{3}
\setlength{\emergencystretch}{2em}
\setlength{\LTpre}{6pt}
\setlength{\LTpost}{6pt}
\setlength{\LTcapwidth}{\textwidth}
\newlength{\TableWidth}
\setlist{itemsep=3pt,topsep=5pt,parsep=0pt,leftmargin=*}
\SetAlgoNoLine
\SetAlgoLongEnd
\DontPrintSemicolon
\SetAlgorithmName{Algorithm}{Algorithm}{List of Algorithms}
\SetKwInput{KwIn}{Input}
\SetKwInput{KwOut}{Output}
\SetKw{KwBreak}{break}
\SetKw{KwRet}{return}
\SetAlFnt{\small}
\SetAlCapFnt{\small\bfseries}
\SetAlCapNameFnt{\small\bfseries}
\SetNlSty{textnormal}{}{\hspace{0.45em}}
\setlength{\textfloatsep}{12pt plus 2pt minus 2pt}
\setlength{\floatsep}{12pt plus 2pt minus 2pt}
\setlength{\intextsep}{12pt plus 2pt minus 2pt}
\raggedbottom
\renewcommand{\topfraction}{0.9}
\renewcommand{\bottomfraction}{0.85}
\renewcommand{\textfraction}{0.08}
\renewcommand{\floatpagefraction}{0.7}
\hyphenation{Coauthor CiteSeer CoraFull GB-LCSS Graph-MLP}


\widowpenalty=10000
\clubpenalty=10000
\hypersetup{pdfauthor={Silei Cao, Kun Wang, Jingyu Zhu, Tianxing Wang, Bing Huang},pdftitle={Granular ball local consensus for graph-free node classification under label noise}}

\begin{document}
\begin{center}{\Large\bfseries Supplementary Material}\par\medskip Granular ball local consensus for graph-free node classification under label noise\end{center}

\appendix
\renewcommand{\thetable}{\Alph{section}.\arabic{table}}
\renewcommand{\theHtable}{\Alph{section}.\arabic{table}}


\FloatBarrier

\setcounter{table}{0}

\section{Results of the main comparisons under all noise settings}\label{sec:A}

Accuracies in this supplement are percentages, reported as mean \(\pm\) sample standard deviation over five paired seeds. Bold marks the highest mean within each dataset and noise setting; in Table~C.1, each validation protocol is compared separately. Equations numbered (1)--(23) refer to the main paper, and equations numbered (H.1)--(H.6) are defined in Appendix~H.

\Needspace{120pt}

\subsection{Cora}\label{sec:A.1}

\input{tables/table_A_1.tex}

\Needspace{120pt}

\subsection{CiteSeer}\label{sec:A.2}

\input{tables/table_A_2.tex}

\Needspace{120pt}

\subsection{CoraFull}\label{sec:A.3}

\input{tables/table_A_3.tex}

\Needspace{120pt}

\subsection{Coauthor-CS}\label{sec:A.4}

\input{tables/table_A_4.tex}

\Needspace{120pt}

\subsection{Amazon-Photo}\label{sec:A.5}

\input{tables/table_A_5.tex}

\FloatBarrier

\setcounter{table}{0}

\section{Paired statistical tests}\label{sec:B}

For each setting, the accuracy of the reference method is subtracted from that of GB-LCSS on each of the five paired seeds. Table~\ref{tab:B.1} reports the mean difference and its 95\% confidence interval in percentage points, the two-sided paired \(t\)-test \(p\)-value after Holm correction within the 25 comparisons against each reference method, and the number of seeds on which GB-LCSS is higher. Confidence intervals are unadjusted. Against MLP+CE, the difference remains significant after correction in 24 of the 25 settings; against JoCoR, GB-LCSS is higher on all five seeds in 22 of the 25 settings and significant after correction in 11.

\input{tables/table_B_1.tex}

\FloatBarrier

\setcounter{table}{0}

\section{Sensitivity to noisy validation labels}\label{sec:C}

This experiment checks whether noisy validation labels affect model selection under uniform 0.4 noise. The training split and training-label noise are kept unchanged. Validation labels are independently corrupted with the same uniform noise rate using a separate random seed, while test labels remain unchanged. GB-LCSS and random size-matched grouping use the target construction in Section~3.3 of the main paper; GSDN and PosteL follow the implementations described in Section~4.1.3 of the main paper and Appendix~F.

\input{tables/table_C_1.tex}

Under noisy validation labels, GB-LCSS leads random size-matched grouping by 13.69, 12.35, and 4.87 percentage points on Cora, CiteSeer, and CoraFull, and leads MLP+CE by 13.51, 11.93, and 4.91 points. Its accuracy changes by at most 0.68 points relative to clean validation. The ranking of the six methods is unchanged on Cora; on CiteSeer, GSDN and JoCoR exchange positions, and on CoraFull, GSDN and PosteL exchange positions. Thus, corrupting the validation labels has a small effect on GB-LCSS accuracy in these three uniform 0.4 comparisons.

\FloatBarrier

\setcounter{table}{0}

\section{Additional checks on training targets}\label{sec:D}

This appendix collects three checks on the target construction that are referred to in Section~4.3.1 of the main paper.

Excluding the current node's vote. Averaged over the five noise settings, excluding the node's own vote raises accuracy by 0.08 to 1.03 percentage points depending on the dataset; the smallest difference in any single setting is \ensuremath{-}0.49 points.

Clean training labels. With \(\rho=0\) and otherwise unchanged settings, the accuracy differences from MLP+CE are \ensuremath{-}0.24, +2.00, \ensuremath{-}0.38, \ensuremath{-}0.87, and \ensuremath{-}1.75 percentage points on Cora, CiteSeer, CoraFull, Coauthor-CS, and Amazon-Photo. The refinement improves CiteSeer and lowers accuracy on the other four datasets. With clean labels, a ball whose training nodes all have the same class leaves their targets unchanged; target adjustments are confined to balls containing several classes.

Ties. A tie for the most-voted class occurs in 123 of the 125 dataset--noise--seed runs, in 911 ball instances, and affects 12,788 of the 259,000 training-node targets (4.94\%). For these targets the empirical vote distribution is retained, so the target does not depend on the class numbering.

\FloatBarrier

\setcounter{table}{0}

\section{Results of graph-based soft-target methods under all noise settings}\label{sec:E}

Table~\ref{tab:E.1} reports the five-seed means and sample standard deviations of GB-LCSS, PosteL, ALS, and SALS on the five main datasets and five noise settings. All four methods use the same data splits, training-label noise realizations, and graph-free inference protocol. Implementation details are given in Appendix~F.

\input{tables/table_E_1.tex}

\FloatBarrier

\setcounter{table}{0}

\section{Implementation details and key settings}\label{sec:F}

This appendix gives the data splits, implementation basis, and parameter settings needed to reproduce the comparisons. For each seed from 0 to 4, the NoisyGL percentage-split routine samples disjoint training, validation, and test subsets of \(\lfloor 0.2N\rfloor\), \(\lfloor 0.1N\rfloor\), and \(\lfloor 0.7N\rfloor\) nodes from the full node set without class stratification; because each size is rounded down separately, up to two nodes may belong to none of the three subsets. All nodes remain available as unlabeled graph context during representation learning and partitioning.

Table~\ref{tab:F.1} distinguishes the authors' code, reimplementations, and mechanism adaptations, and identifies the components retained under the common protocol. All methods use the same splits and noise realizations, and every deployed MLP receives only node attributes. Table~\ref{tab:F.2} lists the method-specific settings and how they were selected.

\input{tables/table_F_1.tex}

Optimizer and MLP settings shared across methods are given in Section~4.1.4 of the main paper.

\input{tables/table_F_2.tex}

Settings marked ``Fixed for this study'' were chosen before the full comparison grid and kept unchanged across datasets. Complete configurations are provided in the supplementary repository.

Table~\ref{tab:F.3} lists the GRACE settings used for the five main datasets. Cora and CiteSeer use their public dataset settings; CoraFull and Coauthor-CS use the DBLP setting, and Amazon-Photo uses the PubMed setting. The Dimensions column gives the hidden/projection dimensions. Paired edge-dropping and feature-masking values correspond to the two augmented views. All configurations use a two-layer GCN encoder with weight decay 1e-5.

\input{tables/table_F_3.tex}

These settings were fixed before the label-noise experiments. GRACE pretraining does not use labels or data-split masks.

\FloatBarrier

\setcounter{table}{0}

\section{Results on additional datasets}\label{sec:G}

The six additional datasets were first evaluated under uniform 0.4 noise. Based on these results, CiteSeerFull was selected for an exploratory follow-up covering the other four noise settings. Table~\ref{tab:G.1} summarizes their statistics; data processing and edge homophily follow Table~2 of the main paper.

\input{tables/table_G_1.tex}

Table~\ref{tab:G.2} compares GB-LCSS with the strongest of the seven noisy-label baselines in Table~3 of the main paper and with random size-matched grouping on the six additional datasets under uniform 0.4 noise. Differences are GB-LCSS accuracy minus the corresponding method's accuracy, in percentage points. ``Higher (seeds)'' gives the number of paired runs, out of five, in which GB-LCSS exceeds the strongest baseline. The same definitions apply to Table~\ref{tab:G.3}.

\input{tables/table_G_2.tex}

GB-LCSS has higher mean accuracy than random size-matched grouping on all six datasets, by 4.33--14.25 percentage points, and exceeds the strongest noisy-label baseline on three. The difference from the strongest baseline ranges from \ensuremath{-}5.24 points on PubMed to +1.46 points on CiteSeerFull. Table~\ref{tab:G.3} reports all five noise settings for CiteSeerFull; GB-LCSS has mean accuracy at least as high as the strongest baseline in four settings, and the strongest baseline is higher by 4.20 points under uniform 0.6.

\input{tables/table_G_3.tex}


\FloatBarrier
\setcounter{table}{0}
\section{Analysis of local-consensus supervision}\label{sec:H}
\setcounter{equation}{0}
\renewcommand{\theequation}{H.\arabic{equation}}
\renewcommand{\theHequation}{H.\arabic{equation}}

\ref*{sec:H.1} proves the local-vote error bound in Proposition~1. \ref*{sec:H.2} derives the squared-error identity of Eq.~(16) for each realized target and gives a sufficient condition for lower expected target error. For the voting bound and expected-error analysis, we fix the graph, the frozen representation, the partition, the training-node set, and the clean labels. Randomness comes from independent symmetric label corruption with rate \(\eta\).

\subsection{Reliability of the local vote}\label{sec:H.1}

Fix a training node \(v\) in ball \(B_j\) with \(s_v=|T_{j,-v}|\geq s_{\min}\) voting nodes, and let \(p^{\mathrm{clean}}_{v,c}\) and \(\Delta_v\) be the clean class proportions and the margin of Eq.~(14) in the main paper. Under symmetric noise, each observed label \(\widetilde y_u\) equals the clean label \(y_u\) with probability \(1-\eta\) and equals any particular incorrect class with probability \(\eta/(C-1)\). Writing \(\gamma=1-C\eta/(C-1)\), the expected observed proportion of class \(c\) among the voting nodes is
\begin{equation}\label{eq:H.1}
\mathbb{E}[p_{j,-v,c}]=(1-\eta)p^{\mathrm{clean}}_{v,c}
+\frac{\eta}{C-1}(1-p^{\mathrm{clean}}_{v,c})
=\gamma p^{\mathrm{clean}}_{v,c}+\frac{\eta}{C-1}.
\end{equation}
The expected vote difference between the current node's true class and any competing class \(c\ne y_v\) is therefore
\begin{equation}\label{eq:H.2}
\mathbb{E}[p_{j,-v,y_v}-p_{j,-v,c}]
=\gamma(p^{\mathrm{clean}}_{v,y_v}-p^{\mathrm{clean}}_{v,c})\geq\gamma\Delta_v.
\end{equation}
Equation~\eqref{eq:H.1} shows that symmetric noise shrinks every clean proportion toward the uniform value \(1/C\) by the factor \(\gamma\), and Eq.~\eqref{eq:H.2} shows that, when \(\gamma>0\), the ordering of the clean proportions is preserved in expectation and the clean margin \(\Delta_v\) is reduced to \(\gamma\Delta_v\). The condition \(\eta<(C-1)/C\) in Proposition~1 is exactly \(\gamma>0\); at \(\eta=(C-1)/C\) the observed labels are uniformly distributed and carry no information.

Fix a competing class \(c\ne y_v\) and define, for each voting node \(u\in T_{j,-v}\), the variable \(Z_{u,c}=\mathbf{1}\{\widetilde y_u=y_v\}-\mathbf{1}\{\widetilde y_u=c\}\), which equals \(+1\) if \(u\) votes for the true class of \(v\), \(-1\) if it votes for \(c\), and 0 otherwise. The \(Z_{u,c}\) are independent across voting nodes because the corruptions are independent, they take values in \([-1,1]\), and their average is \(p_{j,-v,y_v}-p_{j,-v,c}\) with expectation given by Eq.~\eqref{eq:H.2}. Hoeffding's inequality [H1] for independent bounded variables gives
\begin{equation}\label{eq:H.3}
\begin{aligned}
\Pr(p_{j,-v,c}\geq p_{j,-v,y_v})
&\leq\exp\!\left[-\frac{s_v\gamma^2(p^{\mathrm{clean}}_{v,y_v}-p^{\mathrm{clean}}_{v,c})^2}{2}\right]\\
&\leq\exp\!\left(-\frac{s_v\gamma^2\Delta_v^2}{2}\right).
\end{aligned}
\end{equation}
The event in Proposition~1 is the union over the \(C-1\) competing classes of the events \(\{p_{j,-v,c}\geq p_{j,-v,y_v}\}\), so the union bound gives the factor \(C-1\) in Eq.~(15) of the main paper; the minimum with one is taken because a probability cannot exceed one. This completes the proof.

The sharpening step of Eq.~(10) moves probability mass toward the unique most-voted class when there is one and retains the empirical distribution when the highest count is tied. It preserves the set of most-voted classes, so the bound also applies to \(q_{j,-v}\). The margin \(\Delta_v\) is defined relative to the true class of each node and can differ between nodes in the same ball. Proposition~1 applies to nodes with \(\Delta_v>0\); Appendix~I.2 reports their proportions on each dataset.

\subsection{Error of the consensus-weighted target}\label{sec:H.2}

Fix any realization of the observed labels and a training node \(v\) with \(|T_{j,-v}|\geq s_{\min}\). Write \(z=e_{\widetilde y_v}\), \(q=q_{j,-v}\), \(\alpha=\alpha_{j,-v}\), and \(e=e_{y_v}\), so that \(t_v=(1-\alpha)z+\alpha q\). Expanding the squared Euclidean distance between \(t_v\) and \(e\) gives
\begin{equation}\label{eq:H.4}
\begin{aligned}
\lVert t_v-e\rVert_2^2={}&(1-\alpha)^2\lVert z-e\rVert_2^2+\alpha^2\lVert q-e\rVert_2^2\\
&+2\alpha(1-\alpha)\langle z-e,q-e\rangle.
\end{aligned}
\end{equation}
Substituting the polarization identity \(2\langle z-e,q-e\rangle=\lVert z-e\rVert_2^2+\lVert q-e\rVert_2^2-\lVert z-q\rVert_2^2\) and collecting terms yields Eq.~(16) of the main paper. The identity holds for every realization of the observed labels, with \(q\) and \(\alpha\) computed from those votes. For a node with fewer than \(s_{\min}\) votes, Eq.~(12) sets \(t_v=z\), and the target error is unchanged.

Two cases follow directly. If the observed label is correct, then \(z=e\), \(L_v(z)=0\), and Eq.~(16) reduces to \(L_v(t_v)=\alpha^2L_v(q)\). If the observed label is corrupted, then \(z\) and \(e\) are distinct one-hot vectors, so \(L_v(z)=2\), and \(L_v(t_v)<2\) whenever \(\alpha>0\) and \(L_v(q)<2\). Summing over the training set, the change in Brier score from the observed one-hot targets \(Z\) to the GB-LCSS targets \(T\) is
\begin{equation}\label{eq:H.5}
\begin{aligned}
\operatorname{BS}(Z)-\operatorname{BS}(T)={}&
\frac{1}{|\mathcal{V}_L|}\sum_{v\in\mathcal{V}_{\mathrm{err}}}\bigl(2-L_v(t_v)\bigr)\\
&-\frac{1}{|\mathcal{V}_L|}\sum_{v\in\mathcal{V}_{\mathrm{clean}}}L_v(t_v),
\end{aligned}
\end{equation}
where \(\mathcal{V}_{\mathrm{err}}=\{v\in\mathcal{V}_L:\widetilde y_v\ne y_v\}\) is the set of training nodes with corrupted labels and \(\mathcal{V}_{\mathrm{clean}}=\mathcal{V}_L\setminus\mathcal{V}_{\mathrm{err}}\) is its complement. The first sum is the error removed on corrupted labels and the second is the error added on correct labels; Table~I.2 reports both. Note that the Brier score can decrease on a corrupted label even when the most probable class of the target remains the observed (wrong) class, because moving probability mass toward the true class already reduces the squared error.

Finally, we give a condition under which the soft target has lower expected error than the observed label. Hold fixed the observed labels of all other training nodes, in addition to the quantities fixed at the start of this appendix. Because \(v\) does not vote on its own target, \(q\) and \(\alpha\) are then constants, and the remaining randomness comes from corruption of the observed label of \(v\). Under symmetric noise this happens with probability \(\eta\), and a corrupted one-hot label has error 2, so \(\mathbb{E}_v L_v(z)=2\eta\). Taking the expectation of Eq.~(16) over the label of \(v\) and dropping the nonpositive last term gives
\begin{equation}\label{eq:H.6}
\mathbb{E}_v L_v(t_v)\leq(1-\alpha)\,2\eta+\alpha L_v(q).
\end{equation}
Hence, whenever \(\alpha>0\) and \(L_v(q)<2\eta\), the expected error of the soft target is strictly smaller than the expected error \(2\eta\) of the observed label. The condition \(L_v(q)<2\eta\) requires the local consensus to be closer to the clean label than a randomly corrupted one-hot label is on average; it concerns the whole distribution \(q\), whereas Proposition~1 concerns only its most-voted class. Section~4.3.1 and Appendix~I.2 evaluate the realized target errors and their decomposition.

\paragraph{Reference for Appendix H}
[H1] W. Hoeffding, Probability inequalities for sums of bounded random variables, J. Am. Stat. Assoc. 58 (1963) 13--30.

\FloatBarrier
\setcounter{table}{0}
\section{Parameter sensitivity and target-error decomposition}\label{sec:I}

\subsection{Parameter sensitivity and representation controls}\label{sec:I.1}

The representation and parameter comparisons cover the five datasets in Table~2 of the main paper under uniform 0.4 noise, using seeds 0--4 and the protocol in Section~4.1. GRACE-based variants use the frozen representations from the main comparisons. Variants share the default configuration's data split, noisy labels, initialization seed, MLP, optimizer, 200-epoch budget, and checkpoint-selection rule.

For these paired runs, targets were constructed on the CPU and the MLP was trained on the AMD Radeon RX 7700 XT. The software settings were PyTorch 2.9.1+rocm7.2.1 with four host threads and scikit-learn 1.8.0 with six threads for 2-means. The default configuration was rerun in this setup and provides the reference for Table~7 and Table~\ref{tab:I.1}. Its results differ slightly from the main comparisons, which used GPU target construction.

The default parameters are listed in Section~4.1.4. Each sensitivity row changes only the indicated parameter. The threshold \(s_{\min}\) counts votes after excluding the current node. With \(\beta=0\), the empirical class distribution still contributes to the target. Random grouping preserves the default ball sizes; raw \(X\) and \(A^2X\) change the representation used for grouping (Section~4.3.3). Table~\ref{tab:I.1} reports accuracy in percent as mean \(\pm\) sample standard deviation over five paired seeds; bold marks the highest mean in each column.

\input{tables/table_I_1.tex}

The realized ball counts help interpret the budget comparison. With \(B_{\max}=128\), splitting stops at 52--53 balls on Cora and 63--66 on CiteSeer, whereas the other three datasets reach 128 balls. The larger budget therefore allows little additional subdivision on the first two datasets. It improves CoraFull on all five seeds, by \(1.57\pm0.54\) percentage points, and lowers its target Brier score from 0.6321 to 0.6026; Coauthor-CS and Amazon-Photo change by \(-0.17\) and \(-0.46\) points.

Reducing the minimum child-ball size \(n_{\min}\) from 32 to 1 changes mean accuracy by \(-1.65\), \(+0.70\), \(-0.05\), \(-0.31\), and \(-0.47\) points on Cora, CiteSeer, CoraFull, Coauthor-CS, and Amazon-Photo. The default minimum size gives a higher mean on four datasets, with differences ranging from 0.05 to 1.65 points.

The voting threshold controls eligibility for consensus supervision. At \(s_{\min}=10\), the fractions of training nodes with a positive consensus weight are 56.53\%, 58.32\%, 99.13\%, 98.98\%, and 89.96\% in the same dataset order. Cora and CiteSeer have both coverage below 60\% and accuracy losses above 5 points (Table~\ref{tab:I.1}). These measurements show how a higher threshold can substantially reduce the number of nodes receiving local-consensus supervision.

\subsection{Error reduction on training targets}\label{sec:I.2}

Equation~\eqref{eq:H.5} separates the change in target error into the reduction on corrupted labels and the increase on initially correct labels. For each seed, both components are divided by the total number of training nodes; their difference is the net reduction in Brier score. Table~\ref{tab:I.2} reports their means over five seeds for the default configuration. Clean training labels provide the reference for this diagnostic.

\input{tables/table_I_2.tex}

The reduction on corrupted labels exceeds the increase on correct labels in all 25 dataset--seed runs, by a factor of 6.7--18.1 in the five-seed means, and the net reduction ranges from 0.1690 on CoraFull to 0.4877 on Coauthor-CS. Random grouping also reduces the average Brier score but yields lower classification accuracy (Table~\ref{tab:I.1}). Average target error alone therefore does not fully explain downstream performance; the choice of nodes contributing to each consensus also matters.

We also counted the training nodes that satisfy the positive-margin condition of Proposition~1: their true class has more clean votes than each competing class among the other training nodes in their ball. Their proportions are 75.64\%, 65.02\%, 36.13\%, 88.94\%, and 84.76\% on Cora, CiteSeer, CoraFull, Coauthor-CS, and Amazon-Photo. CoraFull has the lowest proportion among the five datasets. Table~\ref{tab:I.2} includes all training nodes, so its error reduction summarizes the targets across the full training set.


\end{document}
